\documentclass[11pt]{ctexart}
\usepackage[a4paper,margin=1in]{geometry}
\usepackage{graphicx}
\usepackage{xcolor}
\usepackage{hyperref}
\definecolor{refgray}{gray}{0.45}
\title{\textbf{市场最新观点汇总}\\[0.4em]{\large 市场主线：结构性供给收缩与需求升级正在重塑多个产业链的定价逻辑，而宏观层面的信贷收缩与通胀风险则构成资产定价的新约束。}}
\author{KC桌面 自动生成}
\date{260614}
\begin{document}
\maketitle
\tableofcontents
\newpage
\section{一页摘要}
\begin{itemize}
  \item 中国电池行业产能管控政策启动供给侧改革，有效供给增速预计从2027年开始放缓，头部企业受益于行业洗牌。
  \item 中国5月信贷增速降至7.7\%历史新低，居民部门持续去杠杆，央行主动收紧流动性以修复利率信号机制。
  \item 全球数据中心现场供电需求因电网瓶颈而爆发，燃气轮机和SOFC等现场电源方案渗透率快速提升。
  \item 美联储风险天平从就业下行重新摆向通胀上行，就业走强与石油溢价叠加，降息预期面临修正。
  \item MLCC行业因AI服务器需求结构升级导致有效产能收缩，预计2028年将出现6\%的供给缺口。
  \item 3.5D封装竞争将于2028年启动，供应链再组织者将获得最大价值增量，封装基板均价创历史新高。
  \item 全球铝市场隐形库存即将耗尽，即便中东复产，2026年下半年仍有近100万吨缺口需要填补。
  \item 市场误读2万亿数据中心投资传闻，实际资金指向网络互联而非数据中心建设，芯片才是真正瓶颈。
  \item 光学内容增长与带宽需求强相关，无论CPO还是NPO架构，每颗GPU所需光学引擎数量正在数量级增长。
\end{itemize}
\section{宏观与利率：信贷收缩与通胀风险重塑政策框架}
\textbf{中国信贷收缩进入结构性新常态，而美国通胀风险重新抬头，两者共同对全球资产定价逻辑构成深远影响。}

\begin{itemize}
  \item 中国5月社融增速降至7.7\%历史新低，新增社融2.03万亿元远低于预期的2.32万亿元，人民币贷款余额增速降至5.5\%。
  \item 居民部门去杠杆已从短期波动演变为长期趋势，5月新增居民贷款为负1410亿元，短期和中长期贷款双双录得负增长。
  \item 企业贷款表面尚可但结构偏虚，票据融资新增5570亿元同比多增近5000亿元，剔除票据后实际企业贷款需求疲弱。
  \item 央行主动收紧流动性以修复利率信号机制，通过缩OMO规模和窗口指导将DR007从1.3\%以下拉回1.4\%以上。
  \item 美国就业市场重新走强，近三个月平均新增就业约18.8万人，显著高于维持失业率稳定所需的约5万人盈亏平衡线。
  \item 5月核心PCE预计环比0.36\%、同比3.4\%，PPI传导超预期，机票价格显示石油相关压力仍在向核心服务传导。
  \item 摩根斯坦利指出，若中东冲突不能短期解决，核心PCE在2026年底前低于3\%的可能性正在快速下降。
  \item 美国原油库存（含战略储备）每天减少约200万桶，实物市场油价仍处高位，金融条件因中东冲突而收紧。
  \item 关税脉冲接近尾声，但石油溢价通过机票等渠道的传导才刚刚开始被市场定价，美联储降息预期可能建立在错误前提上。
\end{itemize}
\begin{figure}[htbp]
\centering
\includegraphics[width=0.88\linewidth]{figures/fig_004.jpg}
\caption{Exhibit 1 - Exhibit 1: The longer oil prices remain elevated, the greater the risk of second-round effects on core inflation.}
\end{figure}
\begin{figure}[htbp]
\centering
\includegraphics[width=0.88\linewidth]{figures/fig_005.jpg}
\caption{Exhibit 2 - Exhibit 2: A broad-based acceleration in payrolls so far this year Payrolls change, 000s}
\end{figure}
\begin{figure}[htbp]
\centering
\includegraphics[width=0.88\linewidth]{figures/fig_006.jpg}
\caption{Exhibit 3 - Exhibit 3: The tariff pass-through may be near completion, but a prolonged conflict adds upside risks PCE cumulative tariff pass-through (pp)}
\end{figure}
\begin{figure}[htbp]
\centering
\includegraphics[width=0.88\linewidth]{figures/fig_007.jpg}
\caption{Exhibit 4 - Exhibit 4: The pace of net wealth gains for households have slowed, and nominal consumption has been overshooting wealth-implied target levels Nominal consumption relative to model-implied target level (Bil \$)}
\end{figure}
\begin{figure}[htbp]
\centering
\includegraphics[width=0.88\linewidth]{figures/fig_008.jpg}
\caption{Exhibit 5 - Exhibit 5: There could be some upside risk to the saving rate, if the cyclical factors in the economy turn less supportive in 2H26}
\end{figure}
\begin{figure}[htbp]
\centering
\includegraphics[width=0.88\linewidth]{figures/fig_009.jpg}
\caption{Exhibit 6 - Exhibit 6: US ending stocks of crude oil continue to move lower US Ending Stocks of Crude Oil and Petroleum Products, Thousands.Barrels}
\end{figure}
\begin{figure}[htbp]
\centering
\includegraphics[width=0.88\linewidth]{figures/fig_010.jpg}
\caption{Exhibit 7 - Exhibit 7: Prices of oil in the physical market remain elevated Crude Oil Spot Price FOB (US\$ per Barrel)}
\end{figure}
\begin{figure}[htbp]
\centering
\includegraphics[width=0.88\linewidth]{figures/fig_011.jpg}
\caption{Exhibit 8 - Exhibit 8: Including the SPR, US oil stocks are falling about 2mn barrels per day}
\end{figure}
\begin{figure}[htbp]
\centering
\includegraphics[width=0.88\linewidth]{figures/fig_012.jpg}
\caption{Exhibit 10 - Exhibit 10: US exports of oil and petroleum products have risen, pushing down net imports of these energy goods}
\end{figure}
\begin{figure}[htbp]
\centering
\includegraphics[width=0.88\linewidth]{figures/fig_013.jpg}
\caption{Exhibit 9 - Exhibit 9: US domestic oil production remains broadly unchanged US Domestic Production of Crude Oil (Thous.Barrels per Day)}
\end{figure}
\begin{figure}[htbp]
\centering
\includegraphics[width=0.88\linewidth]{figures/fig_014.jpg}
\caption{Exhibit 11 - Exhibit 11: US imports of oil have decelerated a little while exports have picked up}
\end{figure}
\begin{figure}[htbp]
\centering
\includegraphics[width=0.88\linewidth]{figures/fig_015.jpg}
\caption{Exhibit 12 - Exhibit 12: Financial conditions have tightened following the conflict in the Middle East MS FCI (Change since February 25, federal funds equivalent, bp)}
\end{figure}
\begin{figure}[htbp]
\centering
\includegraphics[width=0.88\linewidth]{figures/fig_016.jpg}
\caption{Exhibit 13 - Exhibit 13: US effective tariff rate Spot Estimates of Tariff Announcements}
\end{figure}
\begin{figure}[htbp]
\centering
\includegraphics[width=0.88\linewidth]{figures/fig_017.jpg}
\caption{Exhibit 14 - Exhibit 14: US Treasury: Customs and excise deposits from tariffs}
\end{figure}
\begin{figure}[htbp]
\centering
\includegraphics[width=0.88\linewidth]{figures/fig_018.jpg}
\caption{Exhibit 15 - Exhibit 15: Cash withdrawals from the DHS - CBP, as a proxy for tariff refunds Cash Withdrawals from DHS, CBP (Fiscal YTD, Bil \$; 4-wk avg ann)}
\end{figure}
\begin{figure}[htbp]
\centering
\includegraphics[width=0.88\linewidth]{figures/fig_019.jpg}
\caption{Exhibit 16 - Exhibit 16: The effect of incoming data on our US GDP tracking estimate Exhibit 17: GDP nowcasts}
\end{figure}
{\small\color{refgray}\textbf{References:} [R002] NOM：中国信贷收缩的“新常态”正在重塑资产定价逻辑; [R004] 摩根斯坦利：市场真正低估的不是通胀，而是美联储转向加息的路径}

\section{AI/算力：芯片瓶颈与光学内容倍增并行}
\textbf{中国AI发展的真正瓶颈在芯片设计与制造而非数据中心基建，而全球光学内容正进入不可逆的倍增通道，封装架构之争不改变需求确定性。}

\begin{itemize}
  \item 摩根斯坦利指出，市场误读2万亿数据中心投资传闻，该数字来自政策文件中的“吸引投资”而非政府直接投入，资金主要流向网络互联而非数据中心建设。
  \item 中国AI发展的真正瓶颈是芯片设计和制造，缺芯比缺机房更紧迫，芯片供给限制才是算力扩张的核心约束。
  \item 若由三大运营商主导建设2万亿数据中心，反而会造成严重供给过剩，对现有数据中心玩家构成利空。
  \item 在光学领域，当前每颗GPU仅配置2-4个光学引擎，NVL576混合架构升至17个，全光CPO架构达35个，远期Feynman架构可能翻倍至70个。
  \item 光学内容增长的核心驱动力是带宽需求，与封装架构（CPO/NPO）的选择无关，时间争议不改变需求上升趋势。
  \item NPO作为过渡方案良率更高、维护更容易，但长期仍将走向CPO甚至片上光学以解决I/O瓶颈和功耗问题。
  \item 当前CPO面临四大挑战：SoIC良率50-60\%、组装良率20-50\%、激光器散热、以及测试与维修复杂度。
  \item 3.5D封装竞争将于2028年启动，台积电的CoPoS面板级封装与Intel的EMIB-T技术将形成多路线竞争格局。
  \item 4月封装基板出货金额同比增20\%但环比回落，每平米均价同比大涨19\%至122.6万日元历史新高，高端封装供需依然紧张。
  \item Rubin封装出货将在2026年夏季进入全速期，最早6月最晚8月，当前市场预期可能尚未充分计入这一轮放量带来的供应链压力。
\end{itemize}
\begin{figure}[htbp]
\centering
\includegraphics[width=0.88\linewidth]{figures/fig_085.jpg}
\caption{Exhibit 1 - Exhibit 1: Optical Content Increasing, Discussion Is More Around Timing}
\end{figure}
\begin{figure}[htbp]
\centering
\includegraphics[width=0.88\linewidth]{figures/fig_086.jpg}
\caption{Exhibit 2 - Exhibit 2: Our CPO Forecast Was Short of Bullish Expectations, but Still Points to Meaningful Growth Opportunity}
\end{figure}
\begin{figure}[htbp]
\centering
\includegraphics[width=0.88\linewidth]{figures/fig_087.jpg}
\caption{Exhibit 3 - Exhibit 3: Preliminary Expectations from 650 Group Put Larger Opportunity as Scale-Up}
\end{figure}
\begin{figure}[htbp]
\centering
\includegraphics[width=0.88\linewidth]{figures/fig_088.jpg}
\caption{Exhibit 4 - Exhibit 4: Market Missing That Regardless of Timing, Optical Content Increasing}
\end{figure}
\begin{figure}[htbp]
\centering
\includegraphics[width=0.88\linewidth]{figures/fig_089.jpg}
\caption{Exhibit 5 - Exhibit 5: Assumptions and Calculations of Optical Engines (OEs) per Rack Architecture Exhibit 6: 6-7mm Optical Engines of Demand From 59K Switches in Model}
\end{figure}
\begin{figure}[htbp]
\centering
\includegraphics[width=0.88\linewidth]{figures/fig_090.jpg}
\caption{Exhibit 7 - Exhibit 7: CPO Challenges CPO Market Trends \& Challenges 2025 nVIDIA: First 1.6T CPO switch}
\end{figure}
\begin{figure}[htbp]
\centering
\includegraphics[width=0.88\linewidth]{figures/fig_091.jpg}
\caption{Exhibit 8 - Exhibit 8: NPO Is Just a Path Towards Eventual CPO or Interposers Transition to Co-Packaged Optics Gen I: Pluggable Optics Chip-to-Module Links Pluggable Modules}
\end{figure}
{\small\color{refgray}\textbf{References:} [R006] NOM：3.5D封装竞争序幕拉开，真正的赢家不是台积电或英特尔，而是供应链的再组织者; [R009] 摩根斯坦利：市场误读了2万亿数据中心投资，真正的瓶颈在芯片而非基建; [R010] 摩根斯坦利：市场低估的不是CPO时间表，而是光学内容正在进入不可逆的倍增通道}

\section{能源与大宗：铝隐形库存耗尽与石油溢价传导}
\textbf{全球铝市场正经历从隐形短缺到显性短缺的质变，而石油溢价通过核心服务渠道的传导才刚刚开始被市场定价。}

\begin{itemize}
  \item JPM预计2026年全球原铝缺口高达170万吨，第二至第四季度累计隐含库存消耗达230万吨，但显性库存基本持平。
  \item 隐形库存（包括生产商、贸易商和消费者持有的场外库存）目前仅剩约两个月的消耗能力，正在充当市场缓冲垫。
  \item 即便霍尔木兹海峡在6月重启，中东铝厂复产仍需一到两个季度，缺口不会因海峡开放而迅速消失。
  \item 下半年仍有近100万吨缺口需要填补，隐形库存见底后缺口将传导至显性库存（全球约170万吨，中国占75\%）。
  \item 中国社会库存本周出现63万吨周度降幅，为2023年12月以来最大，支撑沪铝价格走高。
  \item 海外需要中国铝产品出口填补缺口，LME价格必须维持高位以保持出口套利窗口打开，这一博弈是下一阶段铝价核心驱动。
  \item 中国资本驱动的海外铝产能扩张正在终结全球铝供应长期纪律，印尼等地产能增长值得关注。
  \item 石油方面，美国原油库存持续下降，实物市场油价仍处高位，但期货价格因谈判预期有所下跌。
  \item 摩根斯坦利指出，石油溢价通过机票等核心服务渠道的传导才刚刚开始，若中东冲突持续，核心PCE年底前难以跌破3\%。
\end{itemize}
\begin{figure}[htbp]
\centering
\includegraphics[width=0.88\linewidth]{figures/fig_073.jpg}
\caption{Figure 1 - Figure 1: Global primary aluminum balance, quarterly}
\end{figure}
\begin{figure}[htbp]
\centering
\includegraphics[width=0.88\linewidth]{figures/fig_074.jpg}
\caption{Figure 2 - Figure 2: JPM estimates for Middle East primary aluminum production, quarterly}
\end{figure}
\begin{figure}[htbp]
\centering
\includegraphics[width=0.88\linewidth]{figures/fig_075.jpg}
\caption{Figure 3 - Figure 3: Regional aluminum ingot premiums}
\end{figure}
\begin{figure}[htbp]
\centering
\includegraphics[width=0.88\linewidth]{figures/fig_076.jpg}
\caption{Figure 4 - Figure 4: Global visible aluminum inventories}
\end{figure}
\begin{figure}[htbp]
\centering
\includegraphics[width=0.88\linewidth]{figures/fig_077.jpg}
\caption{Figure 5 - Figure 5: Chinese unwrought aluminum and aluminum product exports LHS: Thousand mt; RHS: Percent change, yoy}
\end{figure}
\begin{figure}[htbp]
\centering
\includegraphics[width=0.88\linewidth]{figures/fig_078.jpg}
\caption{Figure 6 - Figure 6: Weekly change in China visible aluminum inventory (SHFE + regional warehouses). X-axis = weeks post Chinese New Year (week 0 = week closest to start of CNY), Y-axis = kt of aluminum increase / (decrease)}
\end{figure}
\begin{figure}[htbp]
\centering
\includegraphics[width=0.88\linewidth]{figures/fig_079.jpg}
\caption{Figure 8 - Figure 7: Annualized Chinese quarterly primary aluminum production estimates}
\end{figure}
\begin{figure}[htbp]
\centering
\includegraphics[width=0.88\linewidth]{figures/fig_080.jpg}
\caption{Figure 8 - Figure 8: Total visible China aluminum inventory (SHFE + regional warehouses) Y-axis: Thousand mt; X-axis: Weeks around Chinese New Year (0 = week closest to start of CNY)}
\end{figure}
\begin{figure}[htbp]
\centering
\includegraphics[width=0.88\linewidth]{figures/fig_081.jpg}
\caption{Figure 10 - Figure 10: China aluminum semis \& products export arb proxy (China to rest of Asia)}
\end{figure}
\begin{figure}[htbp]
\centering
\includegraphics[width=0.88\linewidth]{figures/fig_082.jpg}
\caption{Figure 11 - Figure 11: Primary aluminium demand growth forecasts by region Percent change yoy}
\end{figure}
\begin{figure}[htbp]
\centering
\includegraphics[width=0.88\linewidth]{figures/fig_083.jpg}
\caption{Figure 12 - Figure 12: Indonesian primary aluminium production by smelter, quarterly Thousand mt, annualised}
\end{figure}
\begin{figure}[htbp]
\centering
\includegraphics[width=0.88\linewidth]{figures/fig_084.jpg}
\caption{Figure 13 - Figure 13: Global primary aluminum balance and supply and demand growth LHS: Thousand mt; RHS: Percent change, yoy}
\end{figure}
{\small\color{refgray}\textbf{References:} [R004] 摩根斯坦利：市场真正低估的不是通胀，而是美联储转向加息的路径; [R008] JPM：铝市场的真正拐点不是霍尔木兹重启，而是“隐形库存”耗尽的那一刻}

\section{电池与新能源：产能管控启动供给侧改革}
\textbf{中国电池行业产能管控政策正在启动真正的供给侧改革，市场过度关注贸易摩擦而低估了国内政策对行业格局的结构性重塑。}

\begin{itemize}
  \item JPM报告指出，电池产能项目审批权限正从省级政府向更高层级集中，大型项目和跨省投资需更高层级审查，新增产能可能纳入配额管理。
  \item 政策采用“老项目老办法”原则，已获批项目不受影响，实际效果主要体现在2027年下半年之后的新增产能上。
  \item 行业有效供给增速从2027年开始放缓，2028年供需格局改善将更为明显，头部企业将受益于行业洗牌。
  \item 过去许多二线电池厂70\%以上利润来自政府补贴，现地方财政吃紧、绿色贷款增速放缓、A股再融资收紧，小厂融资越来越难。
  \item 2026年6月发布的《对外投资条例》将技术、数据、人员输出纳入监管框架，出海门槛提高。
  \item 市场将电池板块回调归因于中欧贸易摩擦和锂价波动，但JPM认为产能管控才是真正改变行业基本面的力量。
  \item 宁德时代被列为行业首选，中创新航、璞泰来近期被上调评级，逻辑是技术强、资金足、运营好的企业在产能管控时代将获得更多份额。
  \item 行业正从“拼规模”转向“拼技术”，补贴退坡和绿色信贷降温加速了这一转型。
\end{itemize}
\begin{figure}[htbp]
\centering
\includegraphics[width=0.88\linewidth]{figures/fig_001.jpg}
\caption{Figure 1 - Figure 1: China green loans growth (y/y)}
\end{figure}
\begin{figure}[htbp]
\centering
\includegraphics[width=0.88\linewidth]{figures/fig_002.jpg}
\caption{Figure 2 - Figure 2: China EV and ESS battery industry capacity utilization rate (based on designed capacity)}
\end{figure}
\begin{figure}[htbp]
\centering
\includegraphics[width=0.88\linewidth]{figures/fig_003.jpg}
\caption{Figure 3 - Figure 3: China EV and ESS battery industry capacity utilization rate of the top 8 players}
\end{figure}
{\small\color{refgray}\textbf{References:} [R001] JPM：市场低估的不是贸易摩擦，而是中国电池产能管控的结构性红利}

\section{电子元件：MLCC供需拐点与价值迁移}
\textbf{MLCC行业正经历由AI服务器需求驱动的结构性供需拐点，有效产能收缩和产品组合升级将推动行业进入超越2017-2018年超级周期的上行周期。}

\begin{itemize}
  \item JPM预计MLCC行业将从2025年的10\%供给过剩转为2028年的6\%供给短缺，转折力度将超过2017-2018年超级周期。
  \item AI服务器所需MLCC（如1005-47uf超高容值产品）良率仅为常规产品的1/3到1/6，导致有效产能增速远低于名义产能增速。
  \item 服务器MLCC需求预计从2025年约1500亿颗增长到2028年近6300亿颗，是行业增长最大引擎。
  \item 4月日本MLCC出口金额同比增16\%、环比增9.1\%，但国内生产金额同比降10.4\%，一涨一跌反映需求结构变化。
  \item 出口ASP同比上涨5.0\%、环比上涨2.4\%，大中华区ASP环比飙升7.4\%，日本厂商正在向客户交付更贵的产品。
  \item 摩根斯坦利上调2026年全球MLCC出货金额预测至174.3亿美元（同比+18.8\%），核心原因是AI/DC对小型大容量MLCC需求持续扩张。
  \item 村田在小尺寸大容量MLCC上具有优势：DC电压下电容衰减小、高频稳定性好，客户愿为高品质产品付溢价。
  \item 若产能跟不上AI服务器需求爆发，村田可能丢失份额，产能扩张节奏是未来竞争关键变量。
  \item 行业预计ASP增长可达50\%以上，与2017-2018年超级周期类似，但本轮由产品结构升级而非单纯需求拉动。
\end{itemize}
\begin{figure}[htbp]
\centering
\includegraphics[width=0.88\linewidth]{figures/fig_020.jpg}
\caption{Figure 1 - Figure 1: MLCC industry revenue and y-y}
\end{figure}
\begin{figure}[htbp]
\centering
\includegraphics[width=0.88\linewidth]{figures/fig_021.jpg}
\caption{Figure 3 - Figure 3: MLCC industry pricing and y-y}
\end{figure}
\begin{figure}[htbp]
\centering
\includegraphics[width=0.88\linewidth]{figures/fig_022.jpg}
\caption{Figure 5 - Figure 5: MLCC industry revenue vs volume CAGR comparison ASP start to move up upon mix improvement and tightening S-D}
\end{figure}
\begin{figure}[htbp]
\centering
\includegraphics[width=0.88\linewidth]{figures/fig_023.jpg}
\caption{Figure 2 - Figure 2: MLCC industry volume and y-y}
\end{figure}
\begin{figure}[htbp]
\centering
\includegraphics[width=0.88\linewidth]{figures/fig_024.jpg}
\caption{Figure 4 - Figure 4: MLCC Revenue and OPM trend for the top 5 players}
\end{figure}
\begin{figure}[htbp]
\centering
\includegraphics[width=0.88\linewidth]{figures/fig_025.jpg}
\caption{Figure 6 - Figure 6: MLCC content comparison by application}
\end{figure}
\begin{figure}[htbp]
\centering
\includegraphics[width=0.88\linewidth]{figures/fig_026.jpg}
\caption{Figure 7 - Figure 7: MLCC content comparison by application for 2025E}
\end{figure}
\begin{figure}[htbp]
\centering
\includegraphics[width=0.88\linewidth]{figures/fig_027.jpg}
\caption{Figure 9 - Figure 9: MLCC player's UTR trend}
\end{figure}
\begin{figure}[htbp]
\centering
\includegraphics[width=0.88\linewidth]{figures/fig_028.jpg}
\caption{Figure 8 - Figure 8: MLCC industry bottoms-up OPM trend}
\end{figure}
\begin{figure}[htbp]
\centering
\includegraphics[width=0.88\linewidth]{figures/fig_029.jpg}
\caption{Figure 10 - Figure 10: MLCC players' inventory days outstanding}
\end{figure}
\begin{figure}[htbp]
\centering
\includegraphics[width=0.88\linewidth]{figures/fig_030.jpg}
\caption{Figure 11 - Figure 11: Long Term MLCC Supply-Procurement Gap trend}
\end{figure}
\begin{figure}[htbp]
\centering
\includegraphics[width=0.88\linewidth]{figures/fig_031.jpg}
\caption{Figure 12 - Figure 12: MLCC industry bottom-up revenue and y-y trend US\textbackslash{}\$bn}
\end{figure}
\begin{figure}[htbp]
\centering
\includegraphics[width=0.88\linewidth]{figures/fig_032.jpg}
\caption{Figure 14 - Figure 14: MLCC Industry Bottom-up Wallet Share Comparison (Value) \%}
\end{figure}
\begin{figure}[htbp]
\centering
\includegraphics[width=0.88\linewidth]{figures/fig_033.jpg}
\caption{Figure 13 - Figure 13: MLCC industry bottom-up revenue y-y comparison for individual players}
\end{figure}
\begin{figure}[htbp]
\centering
\includegraphics[width=0.88\linewidth]{figures/fig_034.jpg}
\caption{Figure 15 - Figure 15: MLCC top 5 auto MLCC sales mix (\%) and Global auto MLCC unit sales mix (\%) comparison}
\end{figure}
\begin{figure}[htbp]
\centering
\includegraphics[width=0.88\linewidth]{figures/fig_035.jpg}
\caption{Figure 16 - Figure 16: OPM comparison between DRAM, NAND, MLCC and Substrate}
\end{figure}
\begin{figure}[htbp]
\centering
\includegraphics[width=0.88\linewidth]{figures/fig_036.jpg}
\caption{Figure 18 - Figure 18: OPM comparison between MLCC players (2018 vs 2028E)}
\end{figure}
\begin{figure}[htbp]
\centering
\includegraphics[width=0.88\linewidth]{figures/fig_037.jpg}
\caption{Figure 17 - Figure 17: Auto sales CAGR by MLCC players (22-25')}
\end{figure}
\begin{figure}[htbp]
\centering
\includegraphics[width=0.88\linewidth]{figures/fig_038.jpg}
\caption{Figure 19 - Figure 19: MLCC OPM trends by players}
\end{figure}
\begin{figure}[htbp]
\centering
\includegraphics[width=0.88\linewidth]{figures/fig_039.jpg}
\caption{Figure 20 - Figure 20: Past 3M Price Performance – SEMCO vs. MLCC Peers}
\end{figure}
\begin{figure}[htbp]
\centering
\includegraphics[width=0.88\linewidth]{figures/fig_040.jpg}
\caption{Figure 21 - Figure 21: Past 12M Price Performance – SEMCO vs. MLCC Peers}
\end{figure}
\begin{figure}[htbp]
\centering
\includegraphics[width=0.88\linewidth]{figures/fig_041.jpg}
\caption{Figure 22 - Figure 22: FTM P/B for MLCC peers over the past 15 years}
\end{figure}
\begin{figure}[htbp]
\centering
\includegraphics[width=0.88\linewidth]{figures/fig_042.jpg}
\caption{Figure 23 - Figure 23: Average FTM P/B for MLCC peers over the past 15 years}
\end{figure}
\begin{figure}[htbp]
\centering
\includegraphics[width=0.88\linewidth]{figures/fig_043.jpg}
\caption{Figure 24 - Figure 24: MLCC industry combined revenue vs. market cap trend}
\end{figure}
\begin{figure}[htbp]
\centering
\includegraphics[width=0.88\linewidth]{figures/fig_044.jpg}
\caption{Figure 25 - Figure 25: MLCC industry combined OP vs market cap trend Mcap in US\textbackslash{}\$bn, OP in US\textbackslash{}\$mn}
\end{figure}
\begin{figure}[htbp]
\centering
\includegraphics[width=0.88\linewidth]{figures/fig_045.jpg}
\caption{Figure 26 - Figure 26: MLCC industry combined OPM vs market cap trend Mcap in US\textbackslash{}\$bn, \% - RHS}
\end{figure}
\begin{figure}[htbp]
\centering
\includegraphics[width=0.88\linewidth]{figures/fig_046.jpg}
\caption{Figure 27 - Figure 27: Murata market cap vs revenue trend Yen bn, Yen mn}
\end{figure}
\begin{figure}[htbp]
\centering
\includegraphics[width=0.88\linewidth]{figures/fig_047.jpg}
\caption{Figure 28 - Figure 28: Taiyo Yuden market cap vs revenue trend Yen bn, Yen mn}
\end{figure}
\begin{figure}[htbp]
\centering
\includegraphics[width=0.88\linewidth]{figures/fig_048.jpg}
\caption{Figure 29 - Figure 29: SEMCO market cap vs revenue trend}
\end{figure}
\begin{figure}[htbp]
\centering
\includegraphics[width=0.88\linewidth]{figures/fig_049.jpg}
\caption{Figure 30 - Figure 30: Yageo market cap vs revenue trend}
\end{figure}
\begin{figure}[htbp]
\centering
\includegraphics[width=0.88\linewidth]{figures/fig_050.jpg}
\caption{Figure 31 - Figure 31: Murata market cap vs MLCC OP trend}
\end{figure}
\begin{figure}[htbp]
\centering
\includegraphics[width=0.88\linewidth]{figures/fig_051.jpg}
\caption{Figure 33 - Figure 33: SEMCO market cap vs MLCC OP trend}
\end{figure}
\begin{figure}[htbp]
\centering
\includegraphics[width=0.88\linewidth]{figures/fig_052.jpg}
\caption{Figure 32 - Figure 32: Taiyo Yuden market cap vs MLCC OP trend}
\end{figure}
\begin{figure}[htbp]
\centering
\includegraphics[width=0.88\linewidth]{figures/fig_053.jpg}
\caption{Figure 34 - Figure 34: Yageo market cap vs MLCC OP trend}
\end{figure}
\begin{figure}[htbp]
\centering
\includegraphics[width=0.88\linewidth]{figures/fig_054.jpg}
\caption{Figure 35 - Figure 35: Murata market cap vs MLCC OPM trend}
\end{figure}
\begin{figure}[htbp]
\centering
\includegraphics[width=0.88\linewidth]{figures/fig_055.jpg}
\caption{Figure 36 - Figure 36: Taiyo Yuden market cap vs MLCC OPM trend}
\end{figure}
\begin{figure}[htbp]
\centering
\includegraphics[width=0.88\linewidth]{figures/fig_056.jpg}
\caption{Figure 37 - Figure 37: SEMCO market cap vs MLCC OPM trend}
\end{figure}
\begin{figure}[htbp]
\centering
\includegraphics[width=0.88\linewidth]{figures/fig_057.jpg}
\caption{Figure 38 - Figure 38: Yageo market cap vs MLCC OPM trend}
\end{figure}
\begin{figure}[htbp]
\centering
\includegraphics[width=0.88\linewidth]{figures/fig_058.jpg}
\caption{Exhibit 2 - Exhibit 2: MLCC export value and ASP}
\end{figure}
\begin{figure}[htbp]
\centering
\includegraphics[width=0.88\linewidth]{figures/fig_059.jpg}
\caption{Exhibit 3 - Exhibit 3: MLCC export value and ASP (Narita)}
\end{figure}
\begin{figure}[htbp]
\centering
\includegraphics[width=0.88\linewidth]{figures/fig_060.jpg}
\caption{Exhibit 4 - Exhibit 4: MLCC export value and ASP (Osaka)}
\end{figure}
\begin{figure}[htbp]
\centering
\includegraphics[width=0.88\linewidth]{figures/fig_061.jpg}
\caption{Exhibit 5 - Exhibit 5: MLCC export value and ASP (KIX)}
\end{figure}
\begin{figure}[htbp]
\centering
\includegraphics[width=0.88\linewidth]{figures/fig_062.jpg}
\caption{Exhibit 7 - Exhibit 7: Major electric device shipments (annual) Note: \$e =\$ MS estimates Exhibit 8: Murata Mfg: Total production and inventories}
\end{figure}
\begin{figure}[htbp]
\centering
\includegraphics[width=0.88\linewidth]{figures/fig_063.jpg}
\caption{Exhibit 9 - Exhibit 9: Murata Mfg: Total production and inventories}
\end{figure}
\begin{figure}[htbp]
\centering
\includegraphics[width=0.88\linewidth]{figures/fig_064.jpg}
\caption{Exhibit 10 - Exhibit 10: Taiyo Yuden: Capacitor output and consol. inventories}
\end{figure}
\begin{figure}[htbp]
\centering
\includegraphics[width=0.88\linewidth]{figures/fig_065.jpg}
\caption{Exhibit 11 - Exhibit 11: Taiyo Yuden: Capacitor output and consol. inventories}
\end{figure}
\begin{figure}[htbp]
\centering
\includegraphics[width=0.88\linewidth]{figures/fig_066.jpg}
\caption{Exhibit 12 - Exhibit 12: Yageo monthly sales}
\end{figure}
\begin{figure}[htbp]
\centering
\includegraphics[width=0.88\linewidth]{figures/fig_067.jpg}
\caption{Exhibit 13 - Exhibit 13: Walsin monthly sales}
\end{figure}
\begin{figure}[htbp]
\centering
\includegraphics[width=0.88\linewidth]{figures/fig_068.jpg}
\caption{Exhibit 15 - Exhibit 15: Number of MLCCs per vehicle Exhibit 16: MLCC sales}
\end{figure}
\begin{figure}[htbp]
\centering
\includegraphics[width=0.88\linewidth]{figures/fig_069.jpg}
\caption{Exhibit 17 - Exhibit 17: MLCC global market share}
\end{figure}
\begin{figure}[htbp]
\centering
\includegraphics[width=0.88\linewidth]{figures/fig_070.jpg}
\caption{Exhibit 18 - Exhibit 18: MLCC global shipment value and volume}
\end{figure}
\begin{figure}[htbp]
\centering
\includegraphics[width=0.88\linewidth]{figures/fig_071.jpg}
\caption{Exhibit 19 - Exhibit 19: Nippon Chemi-Con \& Nichicon aluminum capacitor sales}
\end{figure}
\begin{figure}[htbp]
\centering
\includegraphics[width=0.88\linewidth]{figures/fig_072.jpg}
\caption{Exhibit 20 - Exhibit 20 Capacitor supply chain ```mermaid graph TD}
\end{figure}
{\small\color{refgray}\textbf{References:} [R005] JPM：MLCC市场低估的不是需求增长，而是有效供给的结构性收缩; [R007] 摩根斯坦利：MLCC市场真正被低估的不是需求总量，而是AI驱动下的价值迁移速度}

\section{现场供电：数据中心电网瓶颈催生新产业链}
\textbf{全球数据中心真正的瓶颈不是GPU供应而是电网接不上，现场供电解决方案正从备用选项转变为主电源结构性替代，市场远未充分定价。}

\begin{itemize}
  \item JPM专家电话会指出，电网互联延迟和电力短缺已取代算力增长本身，成为驱动现场供电解决方案需求的最主要因素。
  \item 美国约60\%的数据中心仍依赖电网供电，但超过20\%已开始使用燃气轮机作为主电源，燃气发动机和SOFC占比快速上升。
  \item 电网审批周期动辄三到五年，而数据中心从决策到投运时间窗口被压缩到18个月以内，时间错配迫使运营商自建发电设施。
  \item JPM预测到2030年全球SOFC需求将达到20-30GW，中国政策支持力度类似当年扶持锂电池和光伏。
  \item SOFC全生命周期成本约0.08-0.15美元/kWh（按8-10年算），随着国产化推进成本将继续下探。
  \item 数据中心看中SOFC的灵活性：低温运行、快速启停，适合做综合电源方案的一部分。
  \item 中国玩家在产能和交付速度上有明显优势，某动力龙头2026年AIDC发动机出货目标3500-4000台，SOFC产能规划激进。
  \item 现场供电产业链不是周期性的资本开支脉冲，而是持续多年的结构性增长机会。
\end{itemize}
{\small\color{refgray}\textbf{References:} [R003] JPM：市场低估的不是AI算力需求，而是数据中心现场供电的结构性缺口}

\section{结语}
综合来看，市场当前对多个产业链的结构性变化定价不足，无论是中国电池产能管控、MLCC有效产能收缩、铝隐形库存耗尽，还是光学内容倍增，均指向供给端或需求端的不可逆转变。与此同时，宏观层面的信贷收缩与通胀风险正在重塑政策框架和资产定价逻辑。投资者需要超越短期市场共识，关注这些结构性力量如何在不同资产类别中逐步显现。
\newpage
\section{报告覆盖清单}
\begin{itemize}
  \item [R001] 电池与新能源：产能管控启动供给侧改革 - JPM：市场低估的不是贸易摩擦，而是中国电池产能管控的结构性红利
  \item [R002] 宏观与利率：信贷收缩与通胀风险重塑政策框架 - NOM：中国信贷收缩的“新常态”正在重塑资产定价逻辑
  \item [R003] 现场供电：数据中心电网瓶颈催生新产业链 - JPM：市场低估的不是AI算力需求，而是数据中心现场供电的结构性缺口
  \item [R004] 宏观与利率：信贷收缩与通胀风险重塑政策框架 / 能源与大宗：铝隐形库存耗尽与石油溢价传导 - 摩根斯坦利：市场真正低估的不是通胀，而是美联储转向加息的路径
  \item [R005] 电子元件：MLCC供需拐点与价值迁移 - JPM：MLCC市场低估的不是需求增长，而是有效供给的结构性收缩
  \item [R006] AI/算力：芯片瓶颈与光学内容倍增并行 - NOM：3.5D封装竞争序幕拉开，真正的赢家不是台积电或英特尔，而是供应链的再组织者
  \item [R007] 电子元件：MLCC供需拐点与价值迁移 - 摩根斯坦利：MLCC市场真正被低估的不是需求总量，而是AI驱动下的价值迁移速度
  \item [R008] 能源与大宗：铝隐形库存耗尽与石油溢价传导 - JPM：铝市场的真正拐点不是霍尔木兹重启，而是“隐形库存”耗尽的那一刻
  \item [R009] AI/算力：芯片瓶颈与光学内容倍增并行 - 摩根斯坦利：市场误读了2万亿数据中心投资，真正的瓶颈在芯片而非基建
  \item [R010] AI/算力：芯片瓶颈与光学内容倍增并行 - 摩根斯坦利：市场低估的不是CPO时间表，而是光学内容正在进入不可逆的倍增通道
\end{itemize}
\section{逐篇报告摘录}
\subsection{[R001] JPM：市场低估的不是贸易摩擦，而是中国电池产能管控的结构性红利}
{\small\color{refgray}归类：电池与新能源：产能管控启动供给侧改革}

JPM：市场低估的不是贸易摩擦，而是中国电池产能管控的结构性红利

过去六周，中国电池供应链相关股票从5月初高点回调了15\%至40\%，而同期沪深300指数仅下跌5\%。市场习惯性地将这一轮调整归因于中欧贸易摩擦升级、锂价波动以及资金流向AI科技板块。但JPM最新发布的电池行业深度报告提出了一个根本不同的判断：市场真正忽略的，是中国正在推进的产能管控政策——这并非一次临时性干预，而是电池行业供给侧改革的正式启动。

这份由Rebecca Wen领衔的研报，系统梳理了四条影响电池行业格局的关键线索：国内产能审批收紧、对外投资新规、中欧贸易摩擦以及美国关税与实体清单。其中，国内产能管控被明确标记为“正面影响”，其他三项则被评估为“中性”。这一定位本身就值得产业决策者认真思考——当市场的注意力被贸易摩擦吸引时，真正改变行业基本面的力量正在国内政策端悄然成型。

1. 产能审批权的上收，正在从根本上改变中国电池行业的竞争规则

JPM报告揭示了一个关键的制度变化：电池产能项目的审批权限正在从省级政府向更高层级集中。过去，地方政府在审批电池产能项目时拥有相当程度的自由裁量权，这在一定程度上推动了产能的快速扩张，但也导致了大量低端同质化产能的涌入。新规的核心变化在于，大型电池项目和跨省投资可能需要更高层级的审查与批准，整体产能增量甚至可能纳入配额管理。

这意味着地方政府独立批准新项目的灵活性将大幅降低。报告特别指出，这一政策采用了明确的“老项目老办法”原则——已获批、已完成备案或已开工的项目不受影响。大多数计划在2026年投产的扩产项目已经完成了相关审批程序，因此能够正常推进。政策的实际效果将主要体现在2027年下半年之后的新增产能上。

这一制度设计的精妙之处在于，它没有采取行政命令式的“一刀切”关停，而是通过提高准入门槛来重塑行业生态。对于头部企业而言，近期的扩产计划已基本锁定，过渡期内反而能享受竞争格局改善的红利；而对于中小厂商来说，未来的扩张路径正在被系统性地收窄。

2. 配额制与选择性扩张框架，将加速行业从“规模竞赛”转向“质量竞争”

报告详细描述了新政策框架下的配额分配逻辑：产能利用率、研发投入强度、制造技术水平、环境合规性、碳足迹表现、财务健康度和运营质量，将成为评估扩产申请的关键指标。土地资源、能源配额、扩产许可等政策资源将优先分配给领先制造商。

这实际上是建立了一个基于综合能力的竞争筛选机制。报告引用了JPM此前发布的100页电池周期分析中的一句话——“一朝被蛇咬，十年怕井绳”，指出经过2023至2024年的产能过剩周期后，行业参与者和投资者都变得明显更加谨慎。而中央政府的这一轮政策调整，恰恰是在行业心态已经发生转变的背景下，用制度化的方式锁定了“理性扩张”的方向。

从数据来看，中国电池行业的产能利用率在2022年达到76\%的峰值后，2023至2024年回落至58\%左右。JPM的测算显示，在政策推动下，2025至2026年利用率有望回升至73\%至77\%，但2027年可能再次回落至69\%。值得注意的是，前八大厂商的产能利用率预计在2025至2026年将超过100\%——这组数据清晰地表明，行业分化正在加速，头部企业的产能利用效率远高于行业平均水平。

报告还揭示了另一个被市场低估的变量：政府对电池行业的金融支持正在系统性收缩。央行数据显示，绿色贷款增速从2022年三季度超过40\%的高点持续回落，2024年降至21\%，2025年一季度进一步降至18\%。这一趋势与电池行业的资本开支周期高度吻合。

与此同时，2024年9月出台的《公平竞争审查条例》明确限制了地方政府向特定企业提供优惠补贴、税收减免等选择性支持。报告引用了一个极具说服力的数据：政府补贴占许多二线电池厂商报告净利润的比例超过70\%，而宁德时代这

[... middle omitted ...]

池及材料子行业（如磷酸铁锂正极、隔膜、铜箔）的闭门行业自律协调。

这些协调的核心内容围绕两个关键词展开：一是“维持定价高于成本”，二是“以产能利用率门槛为参照，放缓或暂停新增产能扩张”。报告特别指出，大多数措施并非强制性，效果最终取决于各子行业的供需平衡和集中度。但早期反馈已经显示出积极信号——干法隔膜价格在行业会议后出现回升和企稳，参与者普遍预期定价“不会从这里进一步恶化”，部分甚至预计2026年价格将继续改善。

这一动态值得关注的原因在于，它表明行业正在从“政策驱动”转向“自律驱动”的自我调节阶段。如果这一趋势能够持续，行业利润率的修复可能比市场预期的更为顺畅。

5. 报告尚未完全回答的关键问题：海外扩张的合规成本与政策博弈的边界在哪里

尽管报告将中欧贸易摩擦和美国关税的影响评估为“中性”，但其中隐含的不确定性不容忽视。欧盟正在酝酿的贸易措施——包括进口配额、额外关税和反胁迫工具——虽然目前仍处于讨论阶段，但其政策走向将直接影响中国电池企业的海外布局节奏。

\subsection{[R002] NOM：中国信贷收缩的“新常态”正在重塑资产定价逻辑}
{\small\color{refgray}归类：宏观与利率：信贷收缩与通胀风险重塑政策框架}

这份NOM研报的核心判断并不复杂，但它的冲击力在于数据背后的结构性含义：中国5月信贷增速降至7.7\%，创历史新低，且这种放缓并非短期波动，而是需求端结构性收缩与供给端政策主动调控共同作用的结果。市场如果只将其解读为“经济复苏乏力”的又一个证据，可能低估了这场信贷收缩对资产定价、企业盈利与政策框架的深远影响。

NOM的数据显示，5月新增社融2.03万亿元，远低于市场预期的2.32万亿元，也低于去年同期的2.29万亿元。人民币贷款余额增速降至5.5\%，同样是历史最低。但真正值得注意的是，这些数字背后隐藏着三个相互关联的结构性变化：居民部门持续去杠杆、企业贷款依赖票据融资“注水”、以及央行主动收紧流动性以修复利率信号机制。

这不是一次简单的周期性放缓。它可能标志着中国信用扩张模式正在从“量宽价低”转向“量稳价韧”，而这一转变对股票、债券和房地产的定价含义，远比表面数据更深远。

NOM报告中最令人不安的数据来自居民贷款。5月新增居民贷款为负1410亿元，而去年同期为正值540亿元。更关键的是，无论是短期贷款还是中长期贷款，双双录得负增长：短期贷款减少840亿元，中长期贷款减少570亿元。这在历史上极为罕见。

这意味着居民部门不仅在减少借贷，还在主动偿还存量债务。居民存款数据同样印证了这一点：5月新增居民存款为负1100亿元，而去年同期为4700亿元。这是自2015年以来首次出现5月居民存款负增长。居民并非没有钱，而是钱正在从银行体系流向其他领域——可能是消费、可能是理财，也可能是偿还负债。

对投资者而言，居民去杠杆意味着什么？它直接抑制了消费类资产的盈利预期，尤其是与可选消费、房地产相关的板块。更深远的影响在于，它改变了中国经济增长的内生动能结构。过去十年，居民加杠杆是支撑房地产和消费升级的核心动力，现在这个动力正在逆转。任何依赖居民信贷扩张的商业模式，都需要重新审视其增长天花板。

企业贷款数据看起来并不差：5月新增企业贷款6400亿元，高于去年同期的5300亿元。但NOM的分析揭示了一个关键扭曲——票据融资规模异常庞大。5月新增票据融资5570亿元，而去年同期仅为750亿元。如果剔除票据融资，企业实际贷款需求可能远低于表面数字。

票据融资通常被视为企业短期流动性管理工具，而非真实投资需求的信号。当企业大量依赖票据而非中长期贷款时，说明它们对未来投资回报缺乏信心，或者银行在信贷额度压力下用票据“凑数”。NOM报告指出，企业中长期贷款5月新增为负200亿元，而去年同期为3300亿元。这是一个鲜明的对比：企业不是在借钱投资，而是在减少长期负债。

这意味着什么？企业部门的资本开支意愿依然低迷。对于工业品、设备制造、工程建设等依赖企业投资的行业，这预示着需求端将持续承压。同时，它也暗示着银行体系面临的资产荒可能比市场想象的更为严重——银行有钱，但找不到优质贷款项目，只能通过票据和债券市场消化资金。

5月信贷数据发布前，市场已经注意到央行在主动收紧流动性。NOM报告详细描述了这一过程：央行通过缩减每日逆回购操作规模、窗口指导国有银行减少同业拆借，将DR007从5月低点1.3\%以下推回至1.4\%以上。10年期国债收益率也从1.7\%左右的低位有所回升。

市场容易将央行的这一操作解读为“货币政策转向收紧”的信号，进而担忧利率上行和流动性紧缩。但NOM的判断更为精细：央行的目标不是为降息或降准腾挪空间，而是修复利率走廊的信号功能，防止债券市场形成自我强化的下跌螺旋，并遏制金融空转风险。

NOM明确表示，维持2026年不降息、不降准的基准判断。这意味着央行当前的操作并非短期对冲，而是对货币政策框架的长期调整。如果这一判断成立，那么市场需要适应一个“流动性充裕但利率不降”的新环境。对于债券投资者而言，这意味着利率下行空

[... middle omitted ...]

基础设施和公共服务，其对经济拉动效应的传导链条较长，且对私人部门投资的溢出效应有限。这意味着信用扩张的“总量”可能企稳，但“结构”将更加依赖公共部门。

对投资者而言，这意味着与政府支出相关的领域（如基建、新能源、公用事业）可能获得相对稳定的需求支撑，而与私人消费和投资相关的领域将继续面临压力。同时，政府债券供给增加也可能对债券市场利率形成上行压力，进一步压缩信用利差。

5月M1同比增长5.5\%，高于4月的5.0\%，也显著高于去年同期的2.3\%。M1通常被视为企业活期存款的代表，其回升往往被解读为企业流动性改善或经营活动回暖。但结合信贷数据看，这一回升可能更多是财政存款减少的统计效应，而非企业主动增加活期存款。

5月财政存款增加7100亿元，低于去年同期的8800亿元。财政存款减少意味着资金从政府账户流向企业或居民账户，这会推高M1。但企业活期存款增加的原因可能是政府支出加快，而非企业自身经营现金流改善。如果企业真实信贷需求依然疲弱，M1回升的可持续性就值得怀疑。

\subsection{[R003] JPM：市场低估的不是AI算力需求，而是数据中心现场供电的结构性缺口}
{\small\color{refgray}归类：现场供电：数据中心电网瓶颈催生新产业链}

JPM：市场低估的不是AI算力需求，而是数据中心现场供电的结构性缺口

数据中心缺电这件事，大多数人还停留在“算力增长快，所以电力需求大”的线性理解上。但JPM近期组织的一场专家电话会，揭示了一个更深刻的判断：全球数据中心真正的瓶颈不是GPU供应，也不是芯片算力，而是电网根本接不上。这个缺口正在推动一个全新的现场供电产业链崛起，其市场空间和利润结构，远未被当前股价充分定价。

这份研报的核心信号是：美国约60\%的数据中心仍依赖电网供电，但超过20\%已经开始使用燃气轮机作为主电源，燃气发动机和固体氧化物燃料电池（SOFC）的占比正在快速上升。这不是备用电源的补充逻辑，而是主电源的结构性替代。当电网的互联延迟和电力短缺直接拖累项目时间表时，数据中心运营商被迫自建发电设施，这正在重塑整个电力供应链的竞争格局。

JPM对潍柴动力和应流股份均维持超配评级，但比评级更值得关注的，是报告中关于SOFC技术路线、供应链壁垒、以及中国企业在全球扩张中真实竞争优势的细节。这些判断如果成立，意味着AIDC电源供应链的估值框架需要被重新定义——不是周期性的资本开支脉冲，而是持续多年的结构性增长。

多数投资者关注AI数据中心时，讨论焦点集中在算力芯片、液冷散热、或者光模块的速率升级。但JPM专家电话会的第一层判断非常直接：电网互联延迟和电力短缺，已经取代了“算力增长”本身，成为驱动现场供电解决方案需求的最主要因素。

这个“所以呢”的推论是：数据中心运营商不再把现场供电视为备用选项，而是将其作为项目能否按时上线的关键路径。在美国，电网审批周期动辄三到五年，而数据中心从决策到投运的时间窗口正在被压缩到18个月以内。这种时间错配意味着，那些能够提供快速部署、可靠输出、且具备规模化供应能力的现场供电方案，将获得显著的定价权和订单优先权。

报告特别指出，中国背景的供应商在这一轮竞争中处于有利位置。原因不仅仅是成本优势，更关键的是交付速度和供应链柔性。韩国的船用发动机厂商虽然也拿到了一些数据中心订单，但JPM专家认为其影响有限——更高的成本、更长的交货周期、以及不够灵活的供应链，限制了它们抓住这一波需求的能力。

对投资者而言，这意味着不应简单地将AIDC电源供应链视为“周期股”或“AI主题股”。电网瓶颈具有结构性特征，不太可能在两三年内被大规模基建投资解决。现场供电的渗透率提升，是一个持续多年的过程。

2. SOFC正在从技术储备变成核心竞争赛道，中国政策支持是关键催化剂

固体氧化物燃料电池（SOFC）并不是一项新技术，但过去十年一直处于“实验室到商业化”的漫长爬坡中。JPM专家电话会给出的新信号是：SOFC已经跨过了商业化门槛，正在成为数据中心现场供电的核心技术选项。

报告引用的数据值得注意：专家预计全球SOFC需求到2030年将达到20-30 GW。作为对比，当前全球SOFC的装机量仍然非常小。这个预测隐含的假设是，SOFC在数据中心领域的渗透率将出现非线性增长。支撑这一判断的逻辑有三层。

第一，SOFC的全生命周期成本已经进入了有竞争力的区间。报告给出的数据是 美元/kWh，假设8-10年的使用寿命。随着产能规模扩Daiwa本地化率提升，成本还有进一步下降的空间。这个成本水平，已经能够与部分地区的电网购电价格竞争。

第二，SOFC的技术特性天然匹配数据中心的需求。低运行温度、快速启停能力、以及燃料灵活性，使得SOFC可以很好地与燃气发动机和储能系统集成，形成完整的现场供电解决方案。数据中心运营商对于供电可靠性的要求极高，SOFC的模块化部署和快速响应能力，正在被客户重新评估。

第三，也是最关键的变量——中国政府的政策支持。报告将当前中国对SOFC的支持，类比于早期对锂电池和光伏产业的扶持。这个类比如果成立，意味着SOFC

[... middle omitted ...]

PPT，而是供应链能力的真实投射

当一家传统发动机企业宣布2026年到2030年SOFC产能从30MW扩张到1GW时，市场通常会有两种反应：要么认为这是炒作，要么认为这不可能实现。但JPM的报告给出了一个不同的判断框架——潍柴的产能规划之所以可信，是因为它建立在已有的供应链、成本管理和技术积累之上。

报告特别强调了潍柴SOFC技术的两个差异化优势：更低的工作温度和更快的响应速度。这两个参数对于数据中心客户而言，直接关系到系统可靠性和部署灵活性。更重要的是，潍柴提供的不是单一产品，而是“燃气发动机+SOFC+储能”的集成解决方案。这种全栈能力，在数据中心运营商眼中意味着更低的系统集成风险和更简单的供应商管理。

另一个常被忽视的细节是潍柴的海外分销策略。报告提到潍柴正在加速建立美国市场的分销网络。对于一家中国制造企业而言，进入美国数据中心供应链的门槛不仅仅是产品竞争力，还包括渠道信任、合规认证、以及售后服务能力。潍柴在这方面的进展，可能是其海外收入增长的重要催化剂。

\subsection{[R004] 摩根斯坦利：市场真正低估的不是通胀，而是美联储转向加息的路径}
{\small\color{refgray}归类：宏观与利率：信贷收缩与通胀风险重塑政策框架 / 能源与大宗：铝隐形库存耗尽与石油溢价传导}

摩根斯坦利：市场真正低估的不是通胀，而是美联储转向加息的路径

当绝大多数市场参与者还在讨论“通胀何时回落、降息何时到来”时，摩根斯坦利最新一期美国经济周报发出了一个信号，其含义远比表面数据复杂：美联储的风险天平正在从“就业下行”重新摆向“通胀上行”，而且这种摆动的速度可能比市场定价所反映的更快。

这不是一个关于数据小幅超预期的讨论。这是一份关于美联储政策框架可能发生结构性转变的报告。去年，美联储因为担忧就业下行而降息75个基点。但今天，摩根斯坦利的经济学家团队在首席美国经济学家Michael Gapen的带领下，提出了一个让市场必须重新审视的命题：我们可能同时面临两个导致美联储无法降息、甚至需要加息的场景——需求拉动型通胀和持续性石油溢价。

这份报告的核心贡献不在于它预测了哪一组数据，而在于它构建了一个清晰的决策框架：当前的数据正在同时向这两个场景靠拢。就业市场在走强，关税脉冲接近尾声，但石油价格通过机票等核心服务渠道的传导才刚刚开始被市场定价。报告中最值得关注的判断是：如果中东冲突不能短期内解决，核心PCE在2026年底前低于3\%的可能性正在快速下降，这意味着2027年快速降息的预期可能建立在错误的前提上。

以下是我们从这份报告中提炼出的五个关键洞察，以及一个报告尚未完全回答的关键问题。

1. 就业市场的重新走强正在改变美联储的风险计算，而市场对此定价不足

摩根斯坦利的分析显示，过去三个月的平均新增就业已回升至约18.8万人，这明显高于去年水平，也高于美联储和市场的普遍预期。更重要的是，这一数字远高于该机构估算的维持失业率稳定所需的约5万人的盈亏平衡线。

这里的关键不是就业数字本身，而是它对美联储政策框架的冲击。一年前，美联储认为就业下行风险大于通胀风险，因此降息75个基点。但今天，就业市场不仅没有恶化，反而在走强。报告的逻辑链条非常清晰：如果就业增长持续保持在10万人以上——远高于5万人的盈亏平衡线——失业率将继续下行，劳动力市场将更接近“需求拉动型通胀”场景。

这意味着什么？美联储的降息选项正在被系统性关闭。市场目前定价的降息预期，可能建立在一个已经过时的就业假设之上。对于资产定价而言，这意味着利率敏感型资产的估值逻辑需要被重新审视。

2. 关税脉冲接近尾声，但石油对核心通胀的传导才刚刚开始被市场低估

报告中最具洞察力的部分之一，是对通胀结构的分层拆解。5月核心CPI环比增长0.21\%，符合预期。关税传导效应接近完成，核心商品价格已经转负。这看上去像是好消息。

但问题出在核心PCE——美联储最关注的通胀指标。根据PPI数据的推算，摩根斯坦利估计5月核心PCE环比增长0.36\%，同比增速达到3.4\%。这个数字显著高于市场预期。背后的推手是服务价格，特别是机票价格的加速上涨。

这里有一个二阶效应值得深入思考：石油对核心通胀的传导，目前仍然集中在机票等少数服务类别，看起来“范围有限”。但报告明确指出，如果冲突不能短期解决，这种传导将从“窄”变“宽”，逐步渗透到更广泛的核心服务中。报告的核心判断是：在没有短期解决方案的情况下，核心PCE在2026年底前低于3\%将变得极其困难，2027年快速降息和通胀快速回落的预期将难以实现。

对于投资者而言，这意味着需要重新评估通胀敏感型资产的配置。市场可能低估了石油价格通过服务渠道对核心通胀的传导强度和持续时间。

3. 消费的“财富效应”已经耗尽，储蓄率可能成为下一个意外变量

摩根斯坦利对消费的分析提供了一个重要的观察框架。该机构的财富模型显示，过去一年名义消费一直高于财富水平所隐含的目标水平。2026年第一季度，这种偏离已经达到990亿美元。

这意味着什么？家庭资产负债表虽然仍然健康，但财富增长正在放缓。消费高于财富隐含水平的状态不可持续。报告

[... middle omitted ...]

一个需要警惕的信号。市场的消费乐观情绪可能忽略了储蓄率回升的潜在压力。

4. 石油库存的物理现实与期货市场的金融预期正在脱节，这种脱节本身就是一个风险信号

报告中的石油追踪数据揭示了一个有趣的分化：期货价格因美伊谈判前景而走低，但现货市场的物理价格仍然高企。美国原油库存（含战略石油储备）正以每天约200万桶的速度下降。

这种期货与现货的脱节，本质上是金融市场对“短期冲突解决”的乐观预期与物理市场“库存正在被消耗”的现实之间的背离。如果冲突解决被推迟，期货价格将被现货价格“拉回”现实。反之，如果冲突快速解决，当前的高现货价格将面临快速回落。

对于投资者而言，这种脱节意味着石油相关资产的定价存在一个巨大的二元期权特征。报告没有直接给出概率判断，但物理库存的持续下降是一个不可忽视的硬约束。市场需要问自己：如果冲突在短期内无法解决，当前的期货定价是否充分反映了库存持续消耗的后果？

5. 报告尚未完全回答的关键问题：劳动力市场的“需求拉动”与“石油供给冲击”能否共存？

\subsection{[R005] JPM：MLCC市场低估的不是需求增长，而是有效供给的结构性收缩}
{\small\color{refgray}归类：电子元件：MLCC供需拐点与价值迁移}

JPM：MLCC市场低估的不是需求增长，而是有效供给的结构性收缩

这份JPM的最新研报做了一件多数投资者尚未完成的事：它不仅给出了一个MLCC（多层陶瓷电容）行业到2028年的供需模型，更在模型背后揭示了一个被忽视的底层逻辑——真正驱动行业进入新一轮上行周期的，不是AI服务器带来的需求增量本身，而是需求结构变化引发的有效供给收缩。当市场还在用“周期品”的框架审视这些被动元件厂商时，供给端的结构性变化正在重塑整个行业的盈利曲线。

报告的核心判断清晰而有力：MLCC行业将从2025年的10\%供给过剩，在2028年转为6\%的供给短缺。这个转折的力度，按报告作者的测算，将超过 年的超级周期。而更值得深思的是，这份判断背后的逻辑链条，指向了一个与市场共识截然不同的方向。

1. 真正改变供给曲线的不是产能扩张，而是产品结构升级带来的“有效产能”损耗

理解这份报告，首先要理解一个关键概念：有效产能。MLCC行业的名义产能过去十年保持了约10\%的年增长，但这份研报明确指出，真正决定市场供需平衡的，是经过生产损耗和良率调整后的有效产能。而后者正在经历一个加速下滑的过程。

原因在于AI服务器和汽车电子所需的MLCC，与消费电子所用的产品有本质区别。AI服务器需要的是大尺寸、高电容、高电压规格的MLCC，例如1005-47uf这类超高容值产品。这些产品的制造工艺要求使用侧隙构造法，其装配良率比普通MLCC低3到6倍。报告估计，这种产能惩罚对名义产能增长的负面冲击，在未来三年将从过去五年的低至中个位数百分比，加速至双位数百分比。

这意味着什么？当市场看到MLCC厂商宣布扩产计划时，实际可用的产能增量可能远低于预期。这不是一个临时性的良率爬坡问题，而是由产品结构永久性升级带来的有效供给收缩。每一颗用于AI服务器的MLCC，其占用的产能资源可能是普通消费级产品的数倍。在需求结构快速向高端迁移的过程中，有效供给的弹性正在被系统性压缩。

2. 服务器需求的结构性跃迁，正在将MLCC从一个“被动跟随”的行业变为“主动定价”的行业

报告中的需求预测数据揭示了这一变化的幅度。以单台设备的MLCC用量来看，通用服务器的用量从2016年的约2500颗增长到2025年的约12000颗，而AI服务器在2025年已经达到每台50000颗，到2030年预计将接近148000颗。仅服务器这一个应用领域，其占MLCC总需求的比重将从2025年的3.2\%跃升至2028年的11.1\%。

JPM预测，2025至2028年间MLCC行业TAM（可寻址市场）的年复合增长率将达到24\%，远超2016至2019年的16\%。更重要的是，这一轮增长的驱动力主要来自ASP（平均销售价格）的上行，而非单纯的数量增长。

这里有一个关键的市场理解偏差。许多人认为MLCC是典型的“量增价跌”的被动元件——随着技术成熟和规模扩大，单价会自然下降。但这份报告揭示的图景恰恰相反：在高端应用领域，由于有效供给受限，价格正在成为主动变量。报告预计，2025至2028年间MLCC行业ASP将累计增长超过50\%，其幅度与2016至2018年的超级周期相当。但这一次的驱动力不是阶段性的备货恐慌，而是需求结构的永久性升级。

如果只看行业历史，MLCC的利润率波动极大。2018年行业OPM（营业利润率）达到36\%的峰值后，在2023年回落至15\%。但这份研报预测，行业OPM将在2028年达到38\%，超越2018年的高点。

这背后的逻辑不是简单的涨价周期，而是产品结构变化带来的利润率重置。高端MLCC的单价更高、技术壁垒更深、客户粘性更强，而能够量产这些产品的供应商极为有限。报告提到，AI服务器级MLCC的供应商只有少数几家能够规模化生产，这本质上构成了一个供给寡头格局。在这种格局下，定价

[... middle omitted ...]

，但报告认为，考虑到行业正处于向AI生态转型的早期阶段，共识盈利预测往往滞后于实际增长。只要增长动能持续，股价动力可能在整个2026年下半年保持强劲。

日本和台湾厂商的估值相对更具吸引力，而韩国同业在经历了更大幅度的上涨后，估值弹性空间已经收窄。这个判断的逻辑是：当行业处于盈利上修的早期阶段时，估值对增长的反应是非线性的，低估值起点往往意味着更大的重估空间。

5. 报告尚未完全回答的关键问题：AI服务器MLCC的良率拐点何时到来？

尽管报告对供需缺口的预测非常清晰，但有一个关键变量尚未被充分讨论：AI服务器MLCC的良率改善路径。报告承认，这类产品的装配良率比普通MLCC低3到6倍，且只有少数供应商能够规模化生产。但良率改善的速度将直接影响有效供给的释放节奏。

如果主要供应商在未来12至18个月内成功将良率提升至接近传统产品的水平，那么有效供给的约束将被大幅缓解，供需缺口可能比预测的更小。反之，如果良率改善慢于预期，供给短缺的程度和持续时间都可能超出当前模型。

\subsection{[R006] NOM：3.5D封装竞争序幕拉开，真正的赢家不是台积电或英特尔，而是供应链的再组织者}
{\small\color{refgray}归类：AI/算力：芯片瓶颈与光学内容倍增并行}

NOM：3.5D封装竞争序幕拉开，真正的赢家不是台积电或英特尔，而是供应链的再组织者

这份NOM研报的核心判断，并非关于某个季度出货量的波动，而是指向一个更根本的结构性变化：半导体先进封装正在从“台积电主导的单一技术路线”转向“多技术路线竞争”格局。报告以4月封装基板出货数据为引子，但真正的洞察在于2028年将启动的3.5D封装量产竞争，以及这一竞争对供应链权力格局的深层含义。

对于关注半导体产业链的投资者而言，当前市场可能过度聚焦于AI芯片本身的算力竞赛，而低估了后端封装技术路线分化所带来的供应链重构机会。NOM这份报告提供了一个关键的观察框架：谁能在3.5D封装时代成为供应链的再组织者，谁就将获得这一轮技术升级中最大的价值增量。

1. 4月封装基板出货回落是暂时扰动，暑期的真正放量才是观察窗口

NOM报告披露的4月数据看似矛盾：出货金额同比增20\%，但环比从3月的273亿日元降至233亿日元。按面积计算，出货量仅同比增1\%，而每平米均价同比大涨19\%至122.6万日元的历史新高。

这一组合信号指向两个关键点：其一，价格持续上行说明高端封装基板的供需依然紧张；其二，量的环比回落可能与Rubin处理器的封装出货节奏有关——NOM判断，由于服务器生产被推迟，Rubin封装出货暂时放缓。报告明确指出，2025年4月正是Blackwell封装进入全速生产的时期，基数效应明显。

更值得关注的是NOM对时间节点的判断：Rubin封装出货将在2026年夏季进入全速期，最早6月，最晚8月。这意味着当前的市场预期可能尚未充分计入这一轮放量带来的供应链压力。对于关注Ibiden等封装基板供应商的投资者，二季度末到三季度初的数据将是验证这一判断的关键窗口。

2. 3.5D封装不是技术升级，而是后端供应链权力格局的重新洗牌

NOM报告最核心的洞察在于对3.5D封装技术路线的分析。3.5D封装不同于传统的2.5D或3D封装，它通过混合键合技术将加速器、HBM和小型推理处理器直接键合在同一封装上。这不仅是技术参数的提升，更是对整个后端供应链组织方式的根本性改变。

当前，台积电通过其3DFabric Alliance平台主导着先进封装的生态。但NOM指出，2028年将启动量产的3.5D封装正在形成两条并行路线：台积电的CoPoS（面板级封装）+玻璃核心基板路线，以及英特尔的EMIB-T技术路线。

这一分化的意义在于：它打破了台积电在先进封装领域的垄断地位，为供应链参与者创造了新的选择空间。对于基板供应商、设备厂商、散热方案提供商而言，不再只有台积电一个客户和标准制定者。这种竞争格局的变化，往往意味着供应链利润分配的重构。

3. 英特尔EMIB-T获得谷歌订单，但产能瓶颈才是真正约束

NOM报告引用了The Information在6月8日的报道：谷歌已向英特尔下单，计划在2028年用EMIB-T技术封装至少300万颗TPU。同时，英伟达也在考虑采用该技术。尽管NOM表示未确认报道的准确性，但这则信息本身已经说明了市场对技术路线分化的预期正在形成。

然而，NOM提出了一个关键的约束条件：Ibiden目前没有能力为EMIB-T技术每年供应300万颗TPU所需的基板。因此，这一消息短期内不会反映在Ibiden的盈利预测中。但报告同时指出，值得关注的是Ibiden是否会宣布产能扩张计划。

这意味着，即便技术路线选择存在不确定性，产能瓶颈本身就是一个值得追踪的投资主题。谁能率先突破产能限制，谁就能在3.5D封装竞争中占据先机。对于投资者而言，关注基板供应商的资本开支计划，可能比关注技术路线选择更具可操作性。

4. 台积电的CoPoS与玻璃基板：时间表清晰，但技术挑战仍未完全解决

NOM报告详细梳理了台积电在年度股东大

[... middle omitted ...]

，这一不确定性可能影响其产品路线图。

5. 从台积电平台迁移到英特尔EMIB-T：这不是简单的替代，而是生态系统的重建

NOM报告点出了一个容易被忽视的技术壁垒：在台积电3DFabric Alliance平台上开发加速器的公司，转向EMIB-T并非易事。原因在于3.5D封装包含了供电和热管理等关键功能，这些功能与后端工艺设计平台深度绑定。

具体而言，供电方面需要与合作伙伴共同开发垂直供电或集成电压调节器（iVRs）；热管理方面需要针对每个芯片优化的微通道冷板。这意味着，任何想要采用英特尔EMIB-T技术的公司，要么依赖英特尔的平台，要么自行建立后端半导体工艺供应链。

这一判断的深层含义是：技术路线的切换不仅是成本问题，更是时间和能力问题。对于谷歌这样拥有足够规模和自研能力的公司，建立自己的后端供应链是可行的。但对于大多数AI芯片设计公司，这种切换可能面临巨大的工程挑战和时间成本。因此，3.5D封装竞争的实际格局，可能比简单的“台积电vs英特尔”二分法更为复杂。

\subsection{[R007] 摩根斯坦利：MLCC市场真正被低估的不是需求总量，而是AI驱动下的价值迁移速度}
{\small\color{refgray}归类：电子元件：MLCC供需拐点与价值迁移}

摩根斯坦利：MLCC市场真正被低估的不是需求总量，而是AI驱动下的价值迁移速度

这份来自摩根斯坦利日本电子元件团队的研报，在4月经济产业省数据公布后，做出了一个大胆的修正。它上调了2026至2028年全球MLCC出货金额预测，幅度之大值得每一位关注半导体供应链和AI硬件投资的读者停下来仔细推敲。

报告的核心判断并非简单的“AI需要更多电容”，而是指向一个更结构性的变化：小型大容量MLCC在AI/数据中心场景中的渗透，正在重塑整个行业的定价逻辑和竞争格局。市场此前可能只看到了量的增长，却低估了价值迁移的速度和深度。

这不仅仅是一份数据更新。它揭示了日本头部MLCC厂商（尤其是村田）如何在看似温和的单价下行趋势中，通过产品组合的跃迁，实现利润率的差异化。对于产业决策者和投资者而言，理解这一层“有量更有价”的叙事，远比跟踪月度出货数字重要。

1. 4月数据揭示了一个关键信号：出口金额增速远超产量，产品组合正在剧烈改善

4月的数据本身存在一个有趣的矛盾。日本国内MLCC产值按美元计算同比下降了10.4\%，但财务省的贸易统计却显示出口金额同比增长了16.0\%。摩根斯坦利的分析师坦承，尚不完全清楚这一差异的原因。但更值得关注的是出口端的内在结构。

出口金额的增长并非由数量驱动。4月MLCC出口数量同比增长10.4\%，金额增速却高出近6个百分点。这意味着平均售价（ASP）在上升。以美元计，出口ASP同比上涨了5.0\%，环比上涨了2.4\%。在大中华区（占4月出口的57.7\%），出口数量仅环比微增0.1\%，但美元出口金额却增长了7.5\%，ASP环比飙升了7.4\%。

这些数字合在一起，指向一个清晰的结论：日本MLCC厂商正在向客户交付更贵的产品。这不是涨价，而是产品组合向高附加值、高单价型号的倾斜。对于长期习惯于ASP每年下降几个百分点的行业而言，这是一个值得警惕的信号——AI服务器和数据中心对小型大容量MLCC的渴求，正在系统性改变出货结构。

2. 摩根斯坦利上调预测的底气：不是周期复苏，是结构性需求跃迁

基于4月数据，摩根斯坦利将2026年全球MLCC出货金额预测从此前的164.8亿美元上调至174.3亿美元，同比增速从12.4\%大幅上调至18.8\%。2027年预测从183.4亿美元上调至208.9亿美元，2028年从203.0亿美元上调至242.5亿美元。

这组数字的含金量在于，它并非简单的线性外推。报告明确指出，驱动增长的核心是“AI/数据中心用高附加值产品需求的持续扩张”。这与 年由疫情驱动的被动元件超级周期有本质区别。那一轮是终端需求爆发带来的全面缺货，而这一次是特定应用场景对特定技术规格的刚需。

新一代GPU的导入，对AI加速器板和芯片用ABF基板的空间提出了极为严苛的要求。要在有限的面积内实现足够的去耦和滤波功能，必须使用小型化、大容量的MLCC。这恰恰是日本厂商（尤其是村田）技术壁垒最高的领域。需求不是普惠的，而是高度选择性的。

3. 村田的护城河：在ASP下行趋势中，用技术溢价重构利润曲线

报告对村田的分析，是理解这轮价值迁移的关键。市场参与者普遍预期同类产品的MLCC单价会温和下降，摩根斯坦利也认同这一判断。但关键在于，村田的竞争策略并非在存量市场打价格战，而是通过技术迭代创造新的增量市场。

报告指出了村田MLCC的三个技术特征：微型化和容量提升的领先地位；施加直流电压后有效电容衰减极小；在高频工作条件下保持电容稳定性。这些特性意味着，当AI加速器在极端功耗和空间约束下运行时，只有少数厂商的产品能够满足设计余量。这为客户支付溢价创造了充分的理由。

因此，投资者不应纠结于“MLCC单价是否在跌”这种笼统的问题。更准确的分析框架是：村田的销售单价（因产品组合升级）可能在上升，而竞争对

[... middle omitted ...]

。如果AI资本开支在 年超出预期，整个供应链的瓶颈可能不在GPU，而在这些看似不起眼的被动元件上。

报告没有详细展开的是：村田之外的二线厂商（如太阳诱电、TDK）能否在这一波需求中分得一杯羹？如果头部厂商产能满载，客户是否会降低对技术指标的苛刻要求，从而为其他厂商打开窗口？这些问题的答案，将决定这轮价值迁移的受益范围是集中于村田一家，还是扩散至整个日本MLCC产业链。

第一，出口ASP的月度变化。如果日本MLCC出口ASP连续3-6个月保持正增长，说明产品组合升级不是一次性事件，而是趋势。尤其要关注大中华区ASP的变动，因为这是AI服务器组装的核心区域。

第二，不同海关关口的ASP差异。报告数据显示，成田机场的MLCC出口ASP高达1.18美分，而关西机场仅为0.25美分。这种巨大差异反映了不同口岸出货的产品档次完全不同。成田主要出货高单价、小批量的高端产品，关西则对应低单价、大批量的通用型号。跟踪这两个关口的出货量和ASP变化，可以提前判断高端产品的放量节奏。

\subsection{[R008] JPM：铝市场的真正拐点不是霍尔木兹重启，而是“隐形库存”耗尽的那一刻}
{\small\color{refgray}归类：能源与大宗：铝隐形库存耗尽与石油溢价传导}

JPM：铝市场的真正拐点不是霍尔木兹重启，而是“隐形库存”耗尽的那一刻

这份JPM最新发布的铝业研报，读起来像一部悬念小说。市场普遍在等待霍尔木兹海峡重启，认为那将是铝价回归理性的转折点。但JPM大宗商品研究团队给出的判断恰恰相反：即便海峡在6月如期开放，中东冶炼厂恢复生产也需要一到两个季度，而在此期间，全球铝市场将经历一个从“隐形短缺”到“显性短缺”的质变过程。

报告的核心主张是：铝市场真正的价格催化剂，不是供应中断本身，而是支撑市场至今的“隐形库存”正在耗尽。当这些藏在生产者、贸易商和消费者手中的库存再也无法填补供需缺口时，中国占全球75\%的显性库存将成为下一个被消耗的对象。而这一过程，将推动LME铝价向4000美元/吨迈进。

这不是一个简单的供需缺口故事。它涉及供应链的传导机制、套利窗口的动态博弈，以及一个正在发生的结构性变化——中国资本驱动的海外铝产能扩张，正在终结全球铝供应的长期纪律。

1. 市场的真实缺口比显性数据大得多，隐形库存正在充当缓冲垫

当前全球铝市场的核心矛盾，存在于两组数据之间。JPM的供需平衡表显示，2026年全球原铝缺口高达170万吨，其中第二季度到第四季度累计隐含库存消耗达230万吨。但与此同时，全球显性铝库存——包括LME、COMEX、SHFE及中国社会库存——基本与2月底水平持平。

这种“数据打架”的背后，是一个被市场低估的缓冲机制：隐形库存。这些库存包括中东生产商存放在亚洲、欧洲和美国的金属，也包括贸易商为避免LME高额仓租而持有的场外库存，以及消费者自身为应对供应链不确定性而保有的缓冲库存。

JPM通过行业会议交流获得的关键反馈是：这些隐形库存目前仅剩约两个月的消耗能力。换句话说，市场正在用“看不见的粮食”度过最艰难的时期，而这份储备即将见底。

这里的含义是：当前铝价的相对稳定本身就是一种假象。市场并没有真正消化170万吨的年度缺口，而是通过消耗供应链各环节的“安全垫”来延缓冲击。一旦这个安全垫耗尽，缺口将被迫向显性库存传导——届时，价格调整的烈度将远高于当前市场预期。

市场对霍尔木兹海峡重启的期待并非没有道理。JPM的基本假设也是海峡将在6月重新开放，中东冶炼厂从第三季度开始恢复生产，并在第四季度和2027年第一季度加速爬坡。

但问题在于时间。JPM的分析显示，即便是有控制的停产，冶炼厂恢复满产也需要一到两个季度。这意味着，即便海峡立即开放，中东地区的实际产量增量在2026年下半年仍然有限。报告中的产量预测图表清晰展示了这一点：中东产量在第二季度触底后，第三季度仅小幅回升，直到第四季度才出现较为明显的环比增长。

这意味着2026年下半年仍有近100万吨的原铝缺口需要填补。而隐形库存的耗尽，将使这部分缺口不得不由显性库存来承担。

对投资者而言，关键的时间窗口是2026年第三季度到第四季度。这是隐形库存耗尽、显性库存开始加速消耗、而中东供应尚未充分恢复的“真空期”。JPM判断，这个窗口将推动LME铝价向4000美元/吨迈进，并将2026年下半年的均价预期上调至3750美元/吨。

3. 中国库存将成为全球铝价的定价锚，套利机制是LME上涨的核心驱动力

当全球市场需要中国来填补中东供应缺口时，一个关键的传导机制开始发挥作用：中国需要维持足够高的铝产品出口，才能弥补海外市场的供应不足。而要让中国持续出口，LME价格必须保持相对于SHFE的溢价，以确保出口套利窗口始终敞开。

JPM的分析指出，中国铝产量虽然可能通过最大化效率提升而短暂突破4500万吨的产能上限，但持续的工业能源使用和排放检查将形成制约，防止供应大幅超标。与此同时，中国国内需求仍在以1.6\%的同比增速增长，如果加上出口拉动的需求，实际增速可达2.5\%。

这意味着，中国自身的铝库存将面临双

[... middle omitted ...]

4. 报告尚未完全回答的关键问题：中国的出口管制将如何改变博弈格局？

JPM在报告中明确提到，中国出于国内供应安全考虑的铝出口管制，是一个显著的看涨风险因素。但报告并未深入展开这一变量的具体情景分析。

这是一个值得追问的问题：如果中国在铝出口上实施类似此前对稀土或镓锗的管制措施，全球铝市场将面临怎样的冲击？当前全球铝市场对中国的依赖度正在上升，中东供应中断已经暴露了供应链的脆弱性。如果中国在此时收紧出口，LME铝价可能远超4000美元/吨的目标价。

但另一方面，中国也有充分的理由避免过度限制出口。铝是中国重要的出口创汇产品，在全球贸易摩擦加剧的背景下，过度限制可能引发贸易伙伴的反制。此外，中国海外投资的铝产能（印度尼西亚、安哥拉、沙特等地）正在快速增长，这些产能虽然短期内无法填补缺口，但长期来看将改变全球铝供应的地理分布。

这个问题的答案，可能取决于中国政府对“国内供应安全”和“出口经济利益”之间的权衡。报告没有给出明确判断，但这恰恰是值得持续跟踪的核心变量。

\subsection{[R009] 摩根斯坦利：市场误读了2万亿数据中心投资，真正的瓶颈在芯片而非基建}
{\small\color{refgray}归类：AI/算力：芯片瓶颈与光学内容倍增并行}

摩根斯坦利：市场误读了2万亿数据中心投资，真正的瓶颈在芯片而非基建

一份未经官方确认的媒体报道，在6月9日引发了市场对国内数据中心和公有云板块的剧烈担忧。Bloomberg报道称，中国政府计划通过发改委在未来五年投入2万亿元人民币建设数据中心，并由电信运营商主导建设和运营。市场的第一反应是：供给过剩、竞争格局恶化、公有云厂商市场份额将被国企侵蚀。

但摩根斯坦利在第一时间发布的这份研报，给出了一个与市场直觉完全不同的判断框架。我们仔细拆解了这份报告的核心逻辑后认为，市场真正需要关注的，不是2万亿数字本身，而是这个数字背后被误读的政策意图，以及中国AI产业链上真正卡脖子的环节。

这份报告的结论值得每一位关注中国科技基础设施投资的读者认真对待：市场低估了政策文件的真实指向，高估了政府主导投资的规模，同时严重忽略了芯片供给才是中国AI发展的真正瓶颈。

这份报告最关键的贡献，是厘清了媒体报道中那个2万亿数字的真实来源。摩根斯坦利团队明确指出，这个数字很可能指向的是2025年1月一份官方政策文件中的表述。该文件提到的是“数据基础设施建设将带动网络、算力和安全设施的建设和升级，预计五年内吸引投资2万亿元”。注意这里的用词是“吸引投资”，而非“政府直接投入”；是“带动网络和算力设施”，而非“数据中心建设”。

更关键的是，这份政策文件所指向的投资方向，是数据中心之间的网络互联，而不是数据中心本身的物理建设。这意味着，2万亿这个数字如果确实存在，它的去向是光纤网络、高速互联、算力调度系统等基础设施，而不是新建大量的数据中心机房。

从产业链角度看，这两者的含义截然不同。数据中心建设投资直接增加机柜供给，可能导致供过于求；而网络互联投资则是提升现有数据中心的利用效率，属于补短板而非扩产能。摩根斯坦利的这一辨析，实际上是在告诉市场：你们担心的供给过剩，可能根本不会发生。

如果仔细阅读摩根斯坦利的分析，会发现一个更深层的矛盾：市场所担心的“政府大规模投资数据中心”，与当前正在执行的数据中心窗口指导政策，在方向上完全相反。

中国的数据中心行业近年来一直受到严格的产能管控。地方政府对新建数据中心项目有明确的能耗指标和审批限制，目的是避免盲目建设和资源浪费。摩根斯坦利在报告中特别提到，当前的政策导向是“平衡数据中心的供需关系”，而非刺激供给。

这种政策自洽性很重要。一个正在通过窗口指导限制数据中心建设的政府，很难同时推出一个2万亿的数据中心建设计划。除非政策方向出现180度大转弯，但报告中没有给出任何支持这一判断的证据。

因此，市场对这份报道的过度反应，本质上是对政策意图的误读。把“网络互联投资”理解为“数据中心建设投资”，把“吸引社会资本”理解为“政府直接投入”，把“提升现有设施效率”理解为“大规模新建产能”——每一个环节的误读都在放大市场的担忧。

摩根斯坦利在报告中做出了一个非常重要的判断，这个判断值得所有关注中国AI产业的读者反复思考：中国AI发展的最大瓶颈是芯片组的生产与设计，而不是数据中心基础设施。

这个判断的逻辑链条很清楚。从供给端看，中国在数据中心建设方面并不缺乏能力。三大电信运营商和第三方IDC服务商在过去几年已经积累了大量的建设和运营经验，国内的数据中心供给能力在全球范围内都处于前列。如果问题出在数据中心不够用，那解决方案是明确的、可行的。

但芯片问题是另一回事。高端AI训练芯片的获取受到出口管制，国产替代芯片在性能和生态上仍有差距。这个问题不是靠花钱就能在短期内解决的。摩根斯坦利暗示，即使数据中心建得再多，没有足够的芯片来填充这些机柜，这些投资也无法转化为有效的算力输出。

这个判断有一个重要的投资含义：如果市场真正理解了这一点，那么关注点应该从数据中心供给转向芯片产业链的突破，以及能够更高效利用现

[... middle omitted ...]

模型能力的领先（阿里的通义千问、字节跳动的Seedance）、以及快速创新的能力。

电信运营商虽然在数据中心物理建设方面有优势，但在上述三个维度上并不具备竞争力。摩根斯坦利的判断很明确：他们不认为这一趋势会逆转。这意味着，即便运营商获得了更多的数据中心资源，也很难将其转化为公有云市场的份额增长。

这个判断对投资人的意义在于：不要因为宏观政策消息就改变对具体公司竞争能力的评估。竞争优势是结构性的，不是靠政策倾斜就能轻易颠覆的。

5. 报告尚未完全回答的关键问题：如果2万亿真的存在，谁来买单？

尽管摩根斯坦利的分析框架很有说服力，但报告本身也留下了一个尚未完全回答的关键问题：如果2万亿的投资计划真的存在，资金从哪里来？

报告提到“吸引投资”而非“政府投入”，但2万亿的规模意味着这不会是单纯的民间资本行为。如果部分资金来自财政或政策性金融工具，那么对电信运营商的资产负债表、资本开支节奏、以及分红能力都会产生实质影响。报告对此着墨不多，这是一个值得进一步深挖的方向。

\subsection{[R010] 摩根斯坦利：市场低估的不是CPO时间表，而是光学内容正在进入不可逆的倍增通道}
{\small\color{refgray}归类：AI/算力：芯片瓶颈与光学内容倍增并行}

摩根斯坦利：市场低估的不是CPO时间表，而是光学内容正在进入不可逆的倍增通道

最近几周，围绕共封装光学（CPO）与近封装光学（NPO）的讨论，让Lumentum、Coherent、Corning等光学标的经历了显著波动。市场似乎将注意力过度集中在“CPO何时量产”这一时间问题上，从而忽视了更根本的结构性事实：无论最终采用哪种架构，每颗GPU所需的光学引擎数量都在以数量级的方式增长。

摩根斯坦利在6月12日发布的最新研报中，给出了一个清晰但被市场低估的核心判断——光学内容的增长与带宽需求强相关，与封装架构的选择无关。这意味着，即便CPO的规模化量产时间表存在不确定性，光学供应链的确定性需求依然在快速上升。对于正在经历回调的这些光学供应商而言，这恰恰是一个需要重新审视的窗口。

这份报告的价值不在于它给出了CPO的精确时间表，而在于它提供了一个更底层的分析框架：把“架构之争”放在一边，先看“带宽需求”这条不可逆的主线。

1. 从每颗GPU 2个光学引擎到35个以上，增长路径已经清晰

摩根斯坦利在报告中给出了一组极具说服力的数据对比。在当前基于可插拔收发器的架构中，每颗GPU仅配置2到4个光学引擎，主要用于scale-out网络。而在NVL576的混合架构中，当CPO被引入用于inter-rack scale-up流量时，每颗GPU的光学引擎附着率跃升至17个左右。若进一步演进到全光学配置，即GPU与网络交换机的连接也转向光子学，这一数字将升至35个。

这还不是终点。报告指出，在未来的Feynman架构和Keyber机架架构中，如果SerDes速度提升至400G，光学引擎数量可维持在35个左右；但如果SerDes停留在200G，这一数字将可能翻倍至70个。换句话说，光学引擎的倍增不是一次性事件，而是随着带宽需求持续增长而不断展开的过程。

这里的核心洞察是：增长路径已经锁定，只是速度问题。市场对CPO时间表的争议，本质上是对“何时”而非“是否”的争议。而“是否”才是决定长期价值的关键变量。

摩根斯坦利明确指出了当前市场的一个认知偏差：投资者在CPO与NPO之间过度纠结，以至于忽略了无论哪种方案胜出，光学内容都在增长这一事实。报告引用Lumentum和Coherent的观点，两家公司均确认光学内容正在进入规模化扩张阶段。

更关键的是，英伟达通过股权投资锁定了Lumentum、Coherent和Corning的产能承诺。这意味着，即便未来架构发生变化，这些供应商的订单也不会被轻易替代。摩根斯坦利认为，这三家供应商仍然是光学内容增长的主要受益者，尤其是在进入scale-up领域后。

这一判断对投资框架的含义是：与其猜测CPO的时间表，不如关注光学供应链的产能扩张节奏和客户锁定程度。那些已经与下游大客户建立深度绑定关系的供应商，在不确定性中拥有更高的确定性。

3. CPO的技术挑战是真实存在的，但也是推动NPO作为过渡方案的原因

摩根斯坦利没有回避CPO面临的现实困难。报告详细列出了四大挑战：第一，CPO将光学器件直接置于电路板上，而光学器件的故障率高于电子器件，导致故障成本上升且无法热插拔；第二，制造成本高昂，CPO的ASP是传统可插拔方案的8到10倍；第三，生态系统成熟度不足，CPO目前集中在英伟达和博通两家，与传统收发器厂商的广泛生态形成鲜明对比；第四，技术复杂度导致良率偏低，尤其是SolC良率仅50-60\%，下游组装良率更低至20-50\%。

正是这些挑战，使得NPO成为当前更务实的选择。NPO将封装与光学分离，降低了生产和组装的难度，同时允许客户在光学组件出现问题时更容易修复。报告引用Tiffany Yeh对WinWay和FOCI的覆盖分析指出，NPO在供应链可复制性和维护便利性上具有明显优势。

[... middle omitted ...]

烈。而CPO将光学引擎直接集成到交换芯片封装中，意味着价值链的重心正在向上游的芯片设计和封装环节转移。英伟达和博通作为CPO的主要推动者，很可能在这一过程中占据利润分配的主导地位。对于Lumentum、Coherent这样的光学组件供应商而言，它们能否在CPO时代维持甚至提升其议价能力，取决于它们在光学引擎设计、封装工艺和客户关系上的不可替代性。

这里仍需验证：当CPO真正进入规模化量产阶段，光学组件供应商的单位利润率是否会因为被整合进更大的芯片封装而受到挤压？还是说，随着光学引擎数量的指数级增长，以量补价的故事能够成立？摩根斯坦利在报告中提到了Lumentum每颗光学引擎约150-200美元的ASP，但并未深入分析这一ASP在CPO成熟后的演变路径。

5. 投资者的观察框架：从“时间表博弈”转向“带宽需求曲线”

基于摩根斯坦利这份报告的逻辑，投资者可能需要调整自己的分析框架。与其在CPO时间表的季度波动中反复博弈，不如将注意力集中在以下几个更具确定性的变量上：

\section{Disclaimer}
{\scriptsize\color{refgray}
This community is a paid learning and information-sharing community created for general educational purposes only. The membership fee is charged solely for access to community discussions, educational content, organization, and operational costs. It is not an advisory fee, management fee, performance fee, brokerage fee, or compensation for personalized financial, investment, legal, tax, or professional advice.\par
I participate and share information solely as an individual and for educational discussion among community members. I am not acting as a financial adviser, investment adviser, broker, dealer, portfolio manager, legal adviser, tax adviser, accountant, fiduciary, or any other licensed professional adviser. Nothing shared in this community constitutes, and should not be interpreted as, financial advice, investment advice, legal advice, tax advice, a recommendation, solicitation, offer, endorsement, instruction, or guarantee to buy, sell, hold, short, trade, or invest in any security, cryptocurrency, fund, derivative, strategy, or other financial product.\par
All content shared in this community, including messages, discussions, links, files, charts, examples, market commentary, watchlists, AI-generated or AI-polished summaries, and third-party materials, is general, impersonal, and educational in nature. It does not take into account any individual's financial situation, investment objectives, risk tolerance, investment horizon, liquidity needs, tax status, legal restrictions, or suitability requirements. No content should be relied upon as suitable for any specific person, account, portfolio, or financial decision.\par
Information shared in this community may be incomplete, inaccurate, outdated, speculative, AI-generated, AI-edited, or based on third-party internet sources that have not been independently verified. I make no representation or warranty as to the accuracy, completeness, timeliness, reliability, or suitability of any information shared.\par
Financial markets involve substantial risk, including the possible loss of principal. Past performance, historical data, hypothetical examples, simulated results, backtests, personal opinions, or educational examples do not guarantee future results. Each member is solely responsible for conducting their own research, exercising independent judgment, and making their own decisions. Before making any financial, investment, legal, or tax decision, members should consult qualified and licensed professionals.\par
Participation in this community, payment of any membership fee, or communication with me or other members does not create any adviser-client, fiduciary, professional, contractual, agency, partnership, employment, or other advisory relationship. By joining or participating in this community, you acknowledge and agree that you are solely responsible for your own decisions, actions, transactions, profits, losses, risks, and consequences, and that I shall not be liable for any direct, indirect, incidental, consequential, financial, legal, tax, or other loss or damage arising from or related to any content, discussion, information, or material shared in this community.\par
}
\end{document}
